% ===========================================================================
\chapter{Verificação, Validação e Considerações Finais}
\label{cap_verificacao_validacao}
% ===========================================================================

A modelagem de casos de uso constitui a base fundamental para a especificação do comportamento dinâmico do sistema, servindo como ponte de comunicação inequívoca entre os especialistas do domínio clínico e a equipe de engenharia de software. Para assegurar que o modelo atenda aos mais elevados padrões de qualidade da indústria e da academia, aplicaram-se métodos formais de verificação estrutural e validação semântica recomendados pela literatura \cite{cockburn2000, iso29148}.

% ===========================================================================
\section{Critérios de Qualidade e Checklist de Verificação}
\label{sec_criterios_qualidade}
% ===========================================================================

A verificação do modelo de casos de uso baseou-se nas diretrizes de Alistair Cockburn \cite{cockburn2000} e nos critérios de consistência da norma ISO/IEC/IEEE 29148:2018 \cite{iso29148}. Cada um dos vinte e dois casos de uso foi submetido ao checklist de conformidade estruturado na Tabela \ref{tab_checklist_qualidade}:

\begin{table}[!ht]
\centering
\caption{Checklist de Verificação da Qualidade dos Casos de Uso}
\label{tab_checklist_qualidade}
\footnotesize
\begin{tabularx}{\textwidth}{p{4.2cm} p{2.2cm} X}
\toprule
\textbf{Critério de Verificação} & \textbf{Referência} & \textbf{Evidência no Modelo do SIGEA-GTT} \\
\midrule
\textbf{Meta Atômica de Valor} & Cockburn (\textit{Sea Level}) & Cada caso de uso atende a um objetivo completo de negócio para o ator, sem fragmentação em cliques pontuais de tela. \\
\textbf{Clareza de Atores e RBAC} & OMG UML 2.5.1 & Segregação inequívoca entre \texttt{ADMIN}, \texttt{PROFESSOR}, \texttt{STUDENT} e o ator não-humano \texttt{SYSTEM\_TIMER}. \\
\textbf{Invariantes e Pré-condições} & ISO/IEC/IEEE 29148 & Todas as pré-condições estabelecem estados estritos de banco (ex.: \texttt{is\_active = true}, \texttt{is\_closed = false}). \\
\textbf{Cobertura de Falhas} & Cockburn (Cenários [e]) & Previsão exaustiva de fluxos de exceção para violações de domínio (@academico), tempo esgotado e notas fora de faixa. \\
\textbf{Semântica de Include/Extend} & OMG UML 2.5.1 & \texttt{<<include>>} restrito a comportamentos obrigatórios e \texttt{<<extend>>} ancorado a pontos de extensão com gatilhos booleanos. \\
\textbf{Rastreabilidade Bidirecional} & ISO/IEC/IEEE 29148 & 100\% dos 25 requisitos funcionais cobertos em matrizes bidirecionais contra os 22 casos de uso. \\
\bottomrule
\end{tabularx}
\fonte{Elaborado pelo autor (2026).}
\end{table}

% ===========================================================================
\section{Validação Clínica Multidisciplinar}
\label{sec_validacao_clinica}
% ===========================================================================

A validação semântica dos casos de uso envolveu a colaboração direta e contínua entre as equipes do Departamento de Computação (DCOMP) e do Departamento de Enfermagem da UFS, sob a orientação do Prof. Dr. Gilton José Ferreira da Silva e da Prof.ª Dr.ª Ana Waleska Silva Patrício.

Os seguintes pontos foram objeto de refinamento multidisciplinar:
\begin{enumerate}
    \item \textbf{Janela Temporal de Auditoria (\texttt{UC12})}: A literatura internacional do IHI \cite{ihi2009} estabelece vinte minutos como tempo ótimo para a revisão primária de prontuários. Confirmou-se a necessidade pedagógica dos alertas sonoros/visuais progressivos (amarelo aos 15 minutos e carmim aos 18 minutos), permitindo que o estudante desenvolva raciocínio clínico sob restrição realista de tempo;
    \item \textbf{Cadeia Causal Integrada (\texttt{UC14} a \texttt{UC17})}: A enfermagem validou que a simples identificação do gatilho clínico é insuficiente para a formação profissional moderna em segurança do paciente. A inclusão articulada do Diagrama de Ishikawa 6M, da Matriz GUT, do Plano 5W2H e do Ciclo PDCA fornece o instrumental metodológico completo preconizado pelas diretrizes da Organização Mundial da Saúde (OMS);
    \item \textbf{Matriz de Confusão e Gabarito (\texttt{UC20})}: A confrontação automatizada das respostas dos alunos com o gabarito docente através do cálculo de Sensibilidade e Concordância NCC MERP reduz a carga de trabalho braçal do professor e transforma o processo de correção em um momento de aprendizado formativo de alto impacto.
\end{enumerate}

% ===========================================================================
\section{Considerações Finais e Próximos Passos}
\label{sec_consideracoes_finais}
% ===========================================================================

O modelo de casos de uso consolidado neste documento atesta a maturidade funcional e comportamental do sistema SIGEA-GTT. Ao unificar o rigor conceitual da Engenharia de Software com as melhores práticas pedagógicas em auditoria hospitalar, estabelece-se um referencial sólido para a construção dos diagramas de interação (Diagramas de Sequência e de Comunicação) e para a homologação das interfaces de usuário e testes automatizados.

Com a modelagem de casos de uso finalizada e rastreada aos vinte e cinco requisitos funcionais, o projeto cumpre integralmente os requisitos de especificação para o Trabalho de Conclusão de Curso em Engenharia de Computação da Universidade Federal de Sergipe.
